\section{Introduction }

This chapter presents the evaluation of AI-Scientist. Following the methodology of Chapter~6, the evaluation is
organized around the qualities the system is designed to deliver --- evidence grounding, claim-level verification,
calibrated confidence, contradiction handling, measured abstention, and traceable reasoning --- and is illustrated
through representative case studies spanning medicine, computer science, biology, and psychology. The chapter then
compares the full multi-agent system against the single-agent and RAG baselines on the shared retrieval substrate,
in line with factuality-oriented evaluation practice \cite{thorne2018fever,wadden2020scifact,min2023factscore}.

The aim of the chapter is to characterize trustworthiness in a richer way, rather than to reduce it to a single
number. It shows \emph{how} the system reasons: which conclusions it grounds, how its confidence responds to
evidence, when it abstains, and how its behaviour differs from the conventional alternatives. Throughout, the
emphasis remains on evidence-aligned, reader-friendly evaluation rather than on a single aggregate score.
This style of reporting is consistent with the move in recent factuality research toward claim-level analysis and
evidence-aware summaries rather than coarse answer-only scoring \cite{min2023factscore,lin2021truthfulqa,es2023ragas}.

\section{Behaviour of the Verification Pipeline }

Across the question set, the full pipeline exhibits a consistent and interpretable pattern. For a research question
within scope, the Orchestrator retrieves a ranked evidence set in which uploaded and curated papers surface ahead
of live-API results, the Claim Extraction agent turns the retrieved abstracts into a small set of focused claims,
and the Verification agent assigns each claim a verdict bound to specific evidence sentences. The executive summary
then describes only what the verified claims support, and the orchestration trace records each step: how many papers
were retrieved, how many claims were extracted, how many critique notes were raised, and whether a second
verification pass was run.

The most important qualitative property is that conclusions are never detached from evidence. A claim marked
supported always carries the sentence that supports it and an overlap score; a claim marked contradicted carries
the opposing sentences; and a claim the evidence cannot settle is reported as insufficient evidence rather than
paraphrased into apparent fact. This is the observable signature of the grounding requirement from Chapter~3.
It also mirrors the evidence-first behaviour that grounding-oriented systems are designed to achieve
\cite{lewis2020rag,guu2020realm,ji2023hallucination}.

\section{Case Study: Medicine }

Consider the medical question of whether metformin improves weight and metabolic outcomes in type~2 diabetes. With
relevant biomedical papers available through the uploaded and live PubMed Central layers, the system extracts
claims about weight management and glycemic outcomes and verifies them against the retrieved abstracts. Claims that
are echoed across more than one independent abstract receive a supported verdict with comparatively high confidence,
because the confidence function rewards corroboration from independent sources. Claims that appear in only a single
abstract are still reported as supported when the evidence aligns, but they attract a Critic note flagging
single-source support and a correspondingly more modest confidence.

This case illustrates two design properties working together. The source-layering behaviour brings genuine
biomedical abstracts into the evidence pool, and the confidence function distinguishes a well-corroborated finding
from one that rests on a single paper. The user therefore sees not just an answer but a graded, source-attributed
assessment, which is exactly the sort of behaviour emphasized in calibrated-answering work
\cite{kadavath2022language,shinn2023reflexion}.

\section{Case Study: Computer Science }

For the question of whether retrieval-augmented generation reduces hallucination in language models, evidence is
drawn from arXiv through the live layer alongside any uploaded papers. The system extracts claims about the effect
of retrieval grounding and verifies them against the retrieved abstracts. This domain is a useful stress test
because the literature contains both supportive findings and explicit caveats. When the retrieved evidence includes
qualifying or negative statements, the contradiction-aware verifier can mark a claim as contradicted or, where the
balance is unclear, as insufficient evidence, and the Critic note records the disagreement. This mirrors recent
concerns about hallucination in large language models and the value of grounded retrieval \cite{ji2023hallucination,huang2023hallucinationsurvey,lewis2020rag}.

The behaviour here demonstrates that the system does not simply confirm the premise of the question. A leading
question phrased to expect a positive answer does not coerce a supported verdict; the verdict follows the evidence,
and where the evidence is mixed the system reports the mixture rather than resolving it in the questioner's favour.

\section{Case Study: Biology }

For the question of whether CRISPR/Cas9 correction of the sickle mutation is a promising treatment for sickle cell
disease, the biomedical live layer supplies abstracts describing gene-editing outcomes. Claims about correction
efficacy and therapeutic promise are verified against these abstracts, and the system again differentiates between
claims that are corroborated across independent sources and claims that rest on a single report. Because the
retrieval layer sorts biomedical results by relevance rather than recency, the evidence pool stays on-topic, which
is visible in the fact that the extracted claims concern the gene-editing question rather than unrelated recent
papers. The biomedical retrieval path benefits from domain-specific embedding and indexing practices similar to
those used in PubMed-oriented retrieval work \cite{gu2021pubmedbert,lee2019biobert}.

This case highlights the value of the live-retrieval implementation details from Chapter~5. Relevance-sorted search
and full-abstract extraction are what make the evidence usable; without them, the verifier would be grounding its
verdicts in off-topic or empty abstracts, a risk that retrieval benchmarks routinely try to control
\cite{petroni2020kilt,thakur2021beir}.

\section{Case Study: Psychology }

For the question of how social-media use relates to adolescent depression and affect, the system retrieves
psychological and clinical abstracts and extracts claims about the direction and strength of the association. This
domain often contains nuanced and partially conflicting findings, and the system's response reflects that nuance:
some claims are supported, some attract single-source or low-confidence critique notes, and overgeneralized claims
--- those phrased in absolute terms --- are deliberately routed to insufficient evidence rather than being
overstated. This is the abstention behaviour operating as intended: the system declines to settle a sweeping claim
that the available abstracts cannot support. That behaviour is particularly appropriate for psychologically nuanced
literature where concise summaries must still remain faithful to the evidence \cite{lin2021truthfulqa,honovich2022true,min2023factscore}.

\section{Source Layering in Practice }

Across all four case studies, the reported source of each retrieved paper confirms that the layered evidence model
behaves as specified. Uploaded papers, when supplied, appear ahead of curated-corpus and live-API papers of
comparable relevance, while genuinely relevant live results still surface in the evidence set. The user interface
makes this visible by labelling each paper as uploaded, curated corpus, or live API. The practical consequence is
that a user can steer the system toward their own provenance by uploading papers, without losing the breadth that
live retrieval provides. Table~\ref{tab:source-layers} summarizes the role of each layer.

\begin{table}[!htbp]
\centering
\small
\begin{tabularx}{0.96\textwidth}{>{\raggedright\arraybackslash}p{0.24\textwidth} >{\raggedright\arraybackslash}p{0.18\textwidth} >{\raggedright\arraybackslash}X}
\toprule
Layer & Priority & Role in the evaluation \\
\midrule
Uploaded papers & Highest & User-provided provenance; surfaces first when relevant; lets the user ground the answer in chosen sources. \\
Curated corpus & Intermediate & Always-available offline grounding; ensures the system works without a network. \\
Live API (PubMed Central, arXiv) & Lowest & Current, broad coverage retrieved at query time; keeps evidence up to date. \\
\bottomrule
\end{tabularx}
\caption{The three evidence layers and their role in verification.}
\label{tab:source-layers}
\end{table}

\section{Confidence Behaviour }

A central evaluation question is whether the per-claim confidence is a usable signal. Observed across the case
studies, confidence behaves monotonically with the factors the design intends: it is highest for claims
corroborated by several independent sources, lower for claims supported by a single source, and lower still --- and
bounded at the floor --- for claims the evidence cannot settle. The explicit penalty for self-only corroboration is
visible in cases where a claim is echoed only by its own source paper: such claims do not receive the high
confidence that genuine multi-source agreement earns. Because confidence is also clamped below an absolute ceiling,
the system never reports certainty, which is the intended posture for a verifier operating on partial evidence.

This behaviour is what makes the confidence number meaningful to a reader. It is not a free-form score emitted by a
model; it is a transparent function of how strong, how corroborated, and how independent the supporting evidence
is, and it moves in the expected direction as those factors change.

\section{Contradiction Handling and Abstention }

The evaluation specifically exercises questions where the literature disagrees or where the evidence is thin. In
these cases the system does not default to a confident answer. When supporting and contradicting evidence coexist
and the counter-evidence is strong and multi-sourced, the verifier returns a contradicted verdict and the report
describes the disagreement. When no sentence aligns sufficiently with a claim, or when a claim is overgeneralized,
the verifier returns insufficient evidence. In both situations the Critic raises a corresponding note so that the
weakness is explicit in the report.

This is the behavioural core of the thesis's evidence-first design. AI-Scientist presents uncertainty and the
literature's disagreement as part of the output, which keeps the final report aligned with the underlying evidence.

\section{The Critic-Driven Re-Verification Loop }

The adaptive loop is observable in the orchestration trace. When the first pass yields a claim that is insufficient,
low-confidence, or single-source, the Critic emits a re-verification decision, the Orchestrator runs a targeted
retrieval for that claim's follow-up query, and the claim is re-verified against the merged evidence. The trace
records this explicitly, for example noting that a particular claim was revisited because its first-pass evidence
was insufficient and that additional papers were checked. Overgeneralized claims are correctly excluded from the
loop, because more evidence cannot rescue a claim that is too broad to settle. The loop therefore concentrates
additional effort precisely where uncertainty was highest, which is the adaptive-reasoning property the thesis
sets out to demonstrate.

\section{Comparison with the Baselines }

Running the same questions through the single-agent and RAG baselines on the shared retrieval substrate isolates
what the multi-agent design adds. The single-agent baseline, which trusts the top retrieved paper, produces a
verdict from a single source and caps its confidence accordingly; it has no mechanism to corroborate across
sources, to detect contradictions, or to abstain on overgeneralized claims beyond a coarse overlap check. The RAG
baseline aggregates supporting overlap across several papers, which is stronger, but it performs no explicit
contradiction analysis and no claim-level decomposition, so it cannot report a contradicted verdict or attribute
support to specific claims, and it likewise caps its confidence.

The full system differs on exactly the axes the thesis emphasizes. It decomposes the answer into claims, binds each
verdict to evidence, computes confidence from corroboration and independence, surfaces contradictions, abstains on
thin or overgeneralized claims, and records a re-verification trace. Table~\ref{tab:baseline-comparison} summarizes
these capability differences, which are properties of the systems' designs rather than tuned outcomes and can be
compared directly with established retrieval and evaluation frameworks \cite{petroni2020kilt,thakur2021beir,es2023ragas}.
The comparison also follows the broader pattern of evaluating retrieval-first systems via explicit capability
deltas, which makes the contribution easier to interpret \cite{ragsurvey2024,lewis2020rag,guu2020realm}.

\begin{table}[!htbp]
\centering
\small
\begin{tabularx}{0.98\textwidth}{>{\raggedright\arraybackslash}X c c c}
\toprule
Capability & Single-agent & RAG & AI-Scientist \\
\midrule
Grounds answer in retrieved evidence & Partial & Yes & Yes \\
Decomposes answer into checkable claims & No & No & Yes \\
Per-claim verdict with attributed evidence & No & No & Yes \\
Explicit contradiction handling & No & No & Yes \\
Confidence from corroboration / independence & No (capped) & No (capped) & Yes \\
Honest abstention on thin / broad claims & Coarse & Coarse & Yes \\
Critic-driven re-verification & No & No & Yes \\
Traceable orchestration record & No & No & Yes \\
\bottomrule
\end{tabularx}
\caption{Capability comparison of the three systems on the shared retrieval substrate.}
\label{tab:baseline-comparison}
\end{table}

\section{Quantitative Results }

The benchmark summary recorded in the project results provides a compact view of the evaluated systems. The full
multi-agent system leads on accuracy and contradiction recall, while also maintaining a strong average confidence
profile, which is consistent with the thesis's emphasis on calibrated, source-aware reasoning over raw fluency
\cite{min2023factscore,es2023ragas,lin2021truthfulqa}.

\begin{table}[!htbp]
\centering
\caption{Core benchmark outcomes from the latest evaluation run}
\begin{tabularx}{0.88\linewidth}{>{\raggedright\arraybackslash}p{0.34\linewidth} >{\raggedright\arraybackslash}X}
\toprule
Metric & Value for the multi-agent system \\
\midrule
Total cases & 12 \\
Accuracy & 0.917 \\
Average confidence & 0.609 \\
Average evidence count & 1.083 \\
Contradiction recall & 0.750 \\
Correct verdicts & 11 of 12 \\
\bottomrule
\end{tabularx}
\end{table}

\begin{figure}[!htbp]
\centering
\includegraphics[draft=false,width=0.9\linewidth]{Fig_7.1.png}
\caption{Comparative performance profile of the evaluated systems}
\label{fig:system_comparison_bar_results}
\end{figure}

\begin{table}[!htbp]
\centering
\caption{Claim-level outcome summary for the multi-agent system}
\begin{tabularx}{0.84\linewidth}{>{\raggedright\arraybackslash}p{0.34\linewidth} >{\raggedright\arraybackslash}X}
\toprule
Outcome & Count \\
\midrule
Supported and correct & 6 \\
Contradicted and correct & 3 \\
Insufficient evidence and correct & 2 \\
Contradicted claims that required tighter coverage & 1 \\
\bottomrule
\end{tabularx}
\end{table}

\begin{figure}[!htbp]
\centering
\begin{tikzpicture}[
  font=\small,
  scale=0.93,
  transform shape,
  box/.style={draw, rounded corners, align=center, text width=30mm, minimum height=10mm, fill=blue!8},
  arrow/.style={-{Stealth[round]}, thick}
]
\node[box] (a) {Claim support};
\node[box, right=6mm of a] (b) {Contradiction handling};
\node[box, right=6mm of b] (c) {Abstention and calibration};
\node[box, right=6mm of c] (d) {Final report quality};

\draw[arrow] (a) -- (b);
\draw[arrow] (b) -- (c);
\draw[arrow] (c) -- (d);
\end{tikzpicture}
\caption{How the multi-agent system strengthens the evaluation chain}
\label{fig:results_flow}
\end{figure}

\section{Domain Generality and Scope Control }

The case studies span four distinct domains served by different live sources, and the pipeline behaves consistently
across all of them: the same retrieve--extract--verify--criticize--synthesize workflow applies whether the question
is biomedical, computational, biological, or psychological. At the same time, the scope guard declines clearly
non-research queries before any analysis begins, returning an explicit message that identifies the non-research
topic. This confirms the domain-generality requirement: the system is broad across legitimate research fields while
still applying its verification design to the task category it is built for.

\section{Result Synthesis }

The evaluation supports the thesis spine in a coherent way.

First, the multi-agent pipeline grounds its conclusions in retrieved evidence and attributes that evidence at the
level of individual claims, which is precisely what the single-agent and RAG baselines do not do.

Second, the per-claim confidence is a meaningful, transparent signal that responds to evidence strength,
corroboration, and source independence, and that never reaches absolute certainty.

Third, the system handles disagreement and limited evidence in a disciplined way, returning contradicted or
insufficient-evidence verdicts and surfacing the corresponding critique.

Fourth, the critic-driven loop demonstrates adaptive reasoning by concentrating additional retrieval and
verification on the specific claims that were weakest after the first pass.

\section{Chapter Conclusion }

This chapter presented the evaluation of AI-Scientist through cross-domain case studies and a comparison with the
single-agent and RAG baselines. The results show that the multi-agent design grounds and attributes its
conclusions, calibrates confidence to evidence, handles contradiction and abstention honestly, and reasons
adaptively through its critic loop --- capabilities the conventional baselines lack on the same evidence. Taken
together, these findings support the central thesis claim: AI-Scientist is best understood as a reproducible,
evidence-grounded, self-critical research assistant that is dependable when the evidence is clear and measured
when it is still emerging. The next chapter discusses the broader meaning of these findings.

